% ============================================================================
%  Entregable 1. Propuesta económica formal (Formulario E-21, FEP01 pp.70 y
%  71). Apartados 1.1, 1.2 y 1.3 del formulario. Las cifras salen de
%  economico/datos/partidas.csv: aquí no se escribe ningún monto a mano.
% ============================================================================

\begin{resumenApertura}
\marcador{Qué propone audIT al mandante, en cuatro o cinco líneas, con el valor
total del proyecto: \textbackslash montoPartida\{proyecto\}}
\begin{recibe}
  \item \marcador{Qué recibe el mandante}
\end{recibe}
\end{resumenApertura}

\section{Propuesta de valor económico}\label{sec:eco-valor}

\marcador{Párrafo de contexto: qué muestra la \tabref{eco-valores}, con el valor
de la fase de implementación, el de la fase de Operación por los 36 meses y el
total del proyecto}

\tablaMontos{Valores totales del proyecto}{eco-valores}{implementacion, operacion}

\marcador{Conclusión que se saca de la tabla}

\subsection{Tipo de cambio, reajustabilidad y vigencia}\label{sec:eco-condiciones}

\marcador{Tipo de cambio referencial y su fecha: los del Formulario E-24, en la
página de parámetros de este documento. Cláusulas de reajustabilidad aplicables
y vigencia de la oferta}

\section{Estructura de pagos detallada}\label{sec:eco-pagos}

\marcador{Párrafo de contexto: la fase de implementación se paga contra los hitos
del Formulario E-25 y la fase de Operación en 36 pagos mensuales}

\subsection{Hitos de pago de la fase de implementación}\label{sec:eco-hitos}

\marcador{Párrafo de contexto: qué muestra la \tabref{hitos-e25}. Los porcentajes
son los del Formulario E-25 y el monto de cada hito sale del valor de la
implementación}

\tablaHitosPago

\marcador{Párrafo de contexto: qué muestra la \tabref{eco-criterios}, el criterio
de aceptación vinculado a cada hito}

\begin{tabla}{Criterio de aceptación vinculado a cada hito de pago}{eco-criterios}{L{12mm} X}
Hito & Criterio de aceptación \\ \midrule
H1 & \marcador{criterio de aceptación} \\
\end{tabla}

\subsection{Pagos mensuales de la fase de Operación}\label{sec:eco-mensual}

\marcador{Servicios incluidos, componentes fijos y variables, métricas que
afectan los pagos variables y periodicidad de facturación. El Formulario E-25
fija 36 pagos, de los meses 21 a 56, con el descuento de las multas por nivel de
servicio}

\section{Resumen ejecutivo de valor}\label{sec:eco-resumen}

\marcador{Párrafo de contexto del resumen ejecutivo de valor}

\subsection{Matriz de servicios y productos incluidos}\label{sec:eco-matriz}

\marcador{Qué incluye la oferta, servicio por servicio}

\subsection{Exclusiones explícitas}\label{sec:eco-exclusiones}

El mandante adquiere el hardware de campo y de red especificado en T-11. Estas adquisiciones se detallan en una hoja informativa fuera del flujo de audIT. El mandante también paga la energía de la sala técnica y de los gabinetes y el combustible del grupo electrógeno. audIT incluye en su oferta la infraestructura on-premise, la instalación interna, la nube, la conectividad, el soporte y la mantención.

\subsection{Supuestos y dependencias comerciales}\label{sec:eco-supuestos}

\marcador{Supuestos y dependencias comerciales}

\subsection{Beneficios económicos para el mandante}\label{sec:eco-beneficios}

La programación de patio tiene un beneficio monetario pendiente de acreditar. El caso base registra CLP 0 y la sensibilidad provisional, CLP 90.000.000 a 36 meses. La verificación de acople y el registro térmico adicional tienen un caso base prudente de CLP 0. La sensibilidad provisional del registro térmico es de CLP 60.000.000 a 36 meses. El beneficio del acople queda pendiente de cuantificación y el kilometraje real se acredita en el criterio 26. El recupero de sobreestadías se estima en CLP 120.700.000 brutos al año bajo el supuesto de recuperar el 50\,\% del monto objetado. La parte retenida por Curimón depende del reparto con los transportistas. La anticipación circadiana considera un beneficio de escenario de CLP 238.320.000 a 36 meses, sujeto a calibración. La valorización del expediente verificable sigue pendiente. Los beneficios del mandante se mantienen separados de los ingresos de audIT. El beneficio neto consolidado queda sujeto a validación.

\subsection{Términos y condiciones comerciales relevantes}\label{sec:eco-terminos}

\marcador{Términos y condiciones comerciales relevantes}
